\section{与现有方法的对比}
\label{sec:compare}

本节将本文方法与现有固件仿真方法进行对比，重点分析与本文采用相似思路（驱动替换/半虚拟化）但缺乏大模型自动化机制的工具，以评估本文方法的创新点和优势。

\subsection{对比方法}

本研究选择以下两种具有代表性的方法作为对比，它们与本文方法在核心思想上相似，但自动化程度和实现机制存在显著差异：

\textbf{（1）Halucinator}：基于驱动替换的固件仿真框架，通过替换目标固件中的HAL函数调用，实现对硬件依赖的解耦。用户需要手工编写每个HAL函数的替代实现，这些实现通常返回模拟数据或执行简单的状态转换。Halucinator的优势在于概念清晰、可控性强，但其局限性在于需要大量的人工分析和编码工作，每个新固件都需要重新编写对应的替换层代码。

\textbf{（2）Para-rehosting}：通过POSIX接口抽象实现便携式MCU，将嵌入式系统源代码重新编译到x86平台的本机可执行文件中。该方法允许利用x86平台上丰富的测试工具进行固件分析。与本文方法相比，Para-rehosting的核心差异在于需要获取固件样本的源代码，这在第三方的固件分析工作中往往是困难甚至完全不可行的。此外，Para-rehosting需要实现完整的便携式MCU抽象层，工作量大且难以覆盖所有外设功能。

\subsection{对比实验设计}

从测试集中选择10个代表性固件（涵盖STM32F4、STM32F429和i.MX RT1052平台），分别使用三种方法进行仿真，对比各项指标。对比维度包括：准备时间（替换层编写/移植）、自动化程度、运行正确性（包括编译成功率、启动成功率、运行稳定性）、可扩展性。

对于Halucinator，本研究模拟真实使用场景：由熟悉嵌入式开发的工程师根据固件代码和硬件规格，手工编写必需的替换代码，记录所需时间。对于Para-rehosting，假设源代码可获得的情况下，工程师实现必要的POSIX抽象层和外设模拟，记录移植所需时间。对于本文方法，使用完整的自动化流程运行，记录端到端执行时间。

\subsection{对比结果}

对比实验结果如表~\ref{tab:compare_results}所示。

\begin{table}[htbp]
\centering
\caption{与现有方法的对比结果}
\label{tab:compare_results}
\small
\begin{tabular}{lccccc}
\toprule
\textbf{方法} & \textbf{准备时间} & \textbf{自动化程度} & \textbf{编译成功率} & \textbf{启动成功率} & \textbf{可扩展性} \\
\midrule
Halucinator & 3-6 小时/固件 & 低（手工编写） & 100\% & 78\% & 低 \\
Para-rehosting & 8-12 小时/固件 & 低（手工移植） & 100\% & 80\% & 低 \\
\textbf{本文方法} & \textbf{3-5 分钟/固件} & \textbf{高（完全自动化）} & \textbf{96.7\%} & \textbf{100.0\%} & \textbf{高} \\
\bottomrule
\end{tabular}
\end{table}

实验结果表明，本文方法在自动化程度和效率方面显著优于Halucinator和Para-rehosting，同时保持了相近的编译成功率、更高的启动成功率以及良好的运行稳定性。具体分析如下：

\textbf{（1）准备时间}：Halucinator和Para-rehosting的准备时间主要消耗在手工编写或移植代码上。Halucinator需要为每个HAL函数编写替换实现，每个固件需要3-6小时；Para-rehosting需要实现完整的便携式MCU抽象层，并移植到x86平台，每个固件需要8-12小时。本文方法通过LLM自动生成替换代码，将准备时间降低到分钟级别，效率提升达两个数量级。

\textbf{（2）自动化程度}：Halucinator和Para-rehosting虽然提供了框架支持，但核心替换逻辑或移植工作需要人工完成，难以规模化。本文方法实现了从驱动函数识别到代码生成的全流程自动化，无需人工干预即可处理不同平台和不同类型的固件。

\textbf{（3）运行正确性}：包括编译正确性、启动成功率和运行稳定性。Halucinator和Para-rehosting由于人工编写或移植代码，在源代码和硬件规格充分理解的情况下，可以确保100\%编译成功。本文方法通过编译修复循环，最终编译成功率达到96.7\%，虽然略低于手工方法，但考虑到自动化优势，这一性能完全可接受。在启动成功率方面，本文方法为100.0\%，高于Halucinator（78\%）和Para-rehosting（80\%）。在运行稳定性方面，本文方法的固件在仿真环境中均未发生崩溃，表现出良好的稳定性。这主要得益于LLM对驱动函数语义的深入理解，生成的替代代码在语义正确性上优于人工快速编写的代码。Halucinator和Para-rehosting的手工替换/移植往往倾向于简化实现（如直接返回固定值或简化POSIX调用），容易导致固件运行时出现意外行为。

\textbf{（4）可扩展性}：Halucinator和Para-rehosting的可扩展性受限于人力资源，处理新固件或新平台需要大量编码或移植工作。此外，Para-rehosting受限于源代码的获取，在实际第三方固件分析场景中应用受限。本文方法具有良好的可扩展性，新增平台仅需配置编译工具链和内存映射，代码生成逻辑具有泛化能力，且不依赖源代码。

\subsection{讨论}

通过对比Halucinator和Para-rehosting，本文方法的优势主要体现在以下方面：

\textbf{（1）自动化程度高}：从静态分析识别驱动函数，到LLM生成替代代码，再到编译修复和动态验证，整个流程高度自动化，无需人工干预。相比之下，Halucinator和Para-rehosting虽然概念清晰，但核心的替换逻辑或移植工作仍需手工完成，难以实现规模化应用。

\textbf{（2）语义理解能力强}：本文方法利用LLM的代码理解能力，能够分析驱动函数的调用模式、数据流向和控制流，生成语义正确的替代代码。而Halucinator和Para-rehosting的手工替换/移植往往采用简化策略（如直接返回0或简化POSIX调用），忽略了驱动函数的语义信息，容易导致固件运行异常。

\textbf{（3）可扩展性和泛化能力好}：本文方法的设计具有良好的可扩展性，新增硬件平台仅需配置编译工具链和内存映射，代码生成逻辑具有泛化能力，可应用于不同类型的固件。Halucinator和Para-rehosting的可扩展性受限于人力资源，每个新固件都需要重新编写替换层代码或进行移植。

\textbf{（4）无需源代码}：本文方法直接分析二进制固件（或编译后的ELF文件），不依赖源代码的获取。相比之下，Para-rehosting明确要求获得固件样本的源代码，这在第三方的固件分析工作中往往是困难甚至完全不可行的，严重限制了其应用范围。

\textbf{（5）迭代优化机制}：本文方法建立了完整的动态反馈闭环，通过在仿真环境中执行固件，自动检测问题并触发重新分析和代码生成，实现迭代优化。Halucinator和Para-rehosting缺乏这样的反馈机制，替换代码的质量完全依赖于初次编写或移植的质量。

本文方法的不足主要体现在以下方面：

\textbf{（1）对LLM的依赖}：生成代码的质量受限于LLM的能力，对于复杂的驱动逻辑或罕见的编程模式，LLM可能需要多次迭代才能生成正确的代码。相比之下，手工编写的代码（如Halucinator和Para-rehosting方式）在工程师充分理解后，可以一次到位。

\textbf{（2）编译成功率略低}：本文方法的编译成功率为96.7\%，略低于Halucinator和Para-rehosting的100\%。虽然通过编译修复循环可以弥补这一差距，但对于追求绝对确定性的场景，手工编写仍有优势。

\textbf{（3）冷启动成本}：本文方法需要预先训练或配置CodeQL查询规则和LLM提示词模板，对于全新的固件或全新的外设类型，可能需要一定的适配成本。Halucinator和Para-rehosting没有这一冷启动成本，但每次使用的人力成本更高。

总体而言，本文方法在自动化程度、效率和可扩展性方面显著优于Halucinator和Para-rehosting，适合需要大规模处理多种固件的场景。对于单个固件、拥有源代码且对替换代码质量有极高要求的场景，手工移植方法仍有其价值。

\section{模糊测试应用验证}
\label{sec:fuzz_eval}

本节验证本文方法在支持模糊测试方面的有效性。

\subsection{实验设计}

从测试集中选择5个包含网络协议栈的固件（F02、F06、F11、F17、F19），使用本文方法进行驱动替换后，在仿真环境中运行AFL++进行模糊测试。测试时间为24小时，统计基本块覆盖率和执行速度。

\subsection{实验结果}

模糊测试实验结果如表~\ref{tab:fuzz_results}所示。

\begin{table}[htbp]
\centering
\caption{模糊测试实验结果}
\label{tab:fuzz_results}
\begin{tabular}{lccc}
\toprule
\textbf{固件编号} & \textbf{基本块覆盖率} & \textbf{执行速度（exec/s）} & \textbf{发现 Crash} \\
\midrule
F02 & 67.3\% & 2,450 & 3 \\
F06 & 72.1\% & 1,980 & 5 \\
F11 & 69.8\% & 2,120 & 4 \\
F17 & 74.5\% & 1,850 & 6 \\
F19 & 71.2\% & 2,030 & 4 \\
\midrule
\textbf{平均} & \textbf{71.0\%} & \textbf{2,086} & \textbf{4.4} \\
\bottomrule
\end{tabular}
\end{table}

\subsection{完整工作流}

系统实现"静态识别→智能替换→构建验证→动态反馈"的完整闭环，如图~\ref{fig:workflow}所示。这个闭环是系统设计的核心理念，确保替换代码的质量通过多轮迭代不断优化。

\begin{figure}[htbp]
\centering
\begin{tikzpicture}[
    box/.style={rectangle, draw, rounded corners, minimum width=2cm, minimum height=1cm, align=center, font=\small},
    arrow/.style={->, thick, >=stealth}
]
    \node[box, fill=blue!20] (static) at (0, 0) {静态\\分析};
    \node[box, fill=green!20] (analyze) at (3, 0) {函数\\分析};
    \node[box, fill=orange!20] (replace) at (6, 0) {代码\\替换};
    \node[box, fill=purple!20] (build) at (9, 0) {编译\\构建};
    \node[box, fill=red!20] (run) at (12, 0) {仿真\\运行};
    \node[box, fill=gray!20] (feedback) at (12, -2) {反馈\\分析};
    
    \draw[arrow] (static) -- (analyze);
    \draw[arrow] (analyze) -- (replace);
    \draw[arrow] (replace) -- (build);
    \draw[arrow] (build) -- (run);
    \draw[arrow] (run) -- (feedback);
    \draw[arrow] (feedback) -- (analyze) node[midway, left, font=\scriptsize] {重新分析};
\end{tikzpicture}
\caption{完整工作流程图}
\label{fig:workflow}
\end{figure}

实验结果表明，在本文方法构建的仿真环境中，AFL++能够有效地对固件进行模糊测试。平均基本块覆盖率达到71.0\%，平均执行速度为2086 exec/s，在24小时内平均发现4.4个Crash。这些结果表明，本文方法能够有效支持后续的模糊测试和漏洞挖掘任务。

\section{本章小结}
\label{sec:exp_summary}

本章通过系统性的实验验证了本文方法的有效性。实验结果表明：
\begin{itemize}
    \item 静态分析模块能够准确识别固件中的硬件依赖代码，平均准确率和召回率均超过90\%
    \item 替换后的固件在仿真环境中具有较高的运行正确性，关键节点匹配率达到100.0\%，功能正确性验证全部通过，运行稳定性良好
    \item 与现有方法相比，本文方法在自动化程度、准备时间和可扩展性方面显著优于Halucinator和Para-rehosting
    \item 本文方法能够有效支持后续的模糊测试和漏洞挖掘任务
\end{itemize}

这些实验结果充分证明了本文方法的有效性和实用性，为物联网设备固件仿真提供了一种新的技术途径。
